\subsection{Bonus Part 18: Neighbour-Search Comparison}
A switchable contact-search implementation was added with \texttt{contact\_search\_method=0} (all-pairs) and \texttt{contact\_search\_method=1} (cell-linked).
Empirical candidate-pair scaling from the measured data is approximately $\mathcal{O}(N^{2.00})$ for all-pairs and $\mathcal{O}(N^{1.46})$ for the cell-linked strategy.
\begin{table}[!t]
\scriptsize
\centering
\caption{Part 18 neighbour-search comparison (serial runs)}
\label{tab:part18-neighbour-search}
\begin{tabular}{rcccccccc}
\toprule
$N$ & $t_{tot}^{AP}$ & $t_{tot}^{CL}$ & $S_{tot}$ & $t_c^{AP}$ & $t_c^{CL}$ & $S_c$ & Cand. factor \\
\midrule
200 & 0.0226 & 0.1932 & 0.12 & 0.0212 & 0.1914 & 0.11 & 80.57 \\
1000 & 0.5019 & 0.2230 & 2.25 & 0.4955 & 0.2164 & 2.29 & 102.69 \\
5000 & 12.5362 & 0.3839 & 32.66 & 12.5061 & 0.3554 & 35.18 & 458.85 \\
\bottomrule
\end{tabular}
\end{table}
where $t_{tot}^{AP}, t_{tot}^{CL}$ are total runtimes; $t_c^{AP}, t_c^{CL}$ are contact-stage runtimes; $S_{tot}, S_c$ are the corresponding cell-linked speedups over all-pairs; and Cand. factor is the candidate-pair ratio $N_{cand}^{AP}/N_{cand}^{CL}$.
Detected-contact counts match between methods for each tested $N$, confirming contact-search correctness while reducing candidate checks.
The measured candidate-pair reduction confirms why the baseline all-pairs method behaves as the expensive $\mathcal{O}(N^2)$ reference algorithm.
For small systems (here $N=200$), the cell-linked method is slower than all-pairs because grid build and traversal overheads dominate the reduced pair-check count.
